% Source-derived chapter 05: Proof of Theorem~\ref{thm:C} and Corollary~\ref{cor:D}
% Original source: anticyclotomic-iwasawa-rational-elliptic-eisenstein
\chapter{Proof of Theorem~\ref{thm:C} and Corollary~\ref{cor:D}}
\label{anticyclotomic-iwasawa-rational-elliptic-eisenstein-chapter-05}
\source{anticyclotomic-iwasawa-rational-elliptic-eisenstein}

\begin{remark}[Source section]
\label{anticyclotomic-iwasawa-rational-elliptic-eisenstein-chapter-05-source}
\source{anticyclotomic-iwasawa-rational-elliptic-eisenstein}
\source{anticyclotomic-iwasawa-rational-elliptic-eisenstein}
This chapter reproduces the corresponding section of the retrieved
LaTeX source, including its definitions, statements, examples, and
proofs. It has not yet been translated into Lean.
\notready
\end{remark}

% The source section title was: Proof of Theorem~\ref{thm:C} and Corollary~\ref{cor:D}
\label{sec:CD}
\source{anticyclotomic-iwasawa-rational-elliptic-eisenstein}

\subsection{Preliminaries}
\label{subsec:HPKS}
\source{anticyclotomic-iwasawa-rational-elliptic-eisenstein}

Let $E$, $p$, and $K$ be as in $\S\ref{sec:Zp-twisted}$, and assume in addition that hypotheses (\ref{eq:intro-Heeg}) and (\ref{eq:intro-disc}) hold. Fix an integral ideal $\mathfrak{N}\subset\mathcal{O}_K$ with $\mathcal{O}_K/\mathfrak{N}=\Z/N\Z$. For each positive integer $m$ prime to $N$, let $K[m]$ be the ring class field of $K$ of conductor $m$, and set
\[
G[m]={\rm Gal}(K[m]/K[1]),\quad\quad\mathcal{G}[m]={\rm Gal}(K[m]/K).
\]
Let also $\mathcal{O}_m=\Z+m\mathcal{O}_K$ be the order of $K$ of conductor $m$.

By the theory of complex multiplication, the cyclic $N$-isogeny between complex CM elliptic curves
\[
\mathbb{C}/\mathcal{O}_K\rightarrow\mathbb{C}/(\mathfrak{N}\cap\mathcal{O}_m)^{-1}
\]
defines a point $x_m\in X_0(N)(K[m])$, and fixing a modular parameterization $\pi:X_0(N)\rightarrow E$ we define the \emph{Heegner point} of conductor $m$ by   
\[
P[m]:=\pi(x_m)\in E(K[m]).
\]
Building on this construction, one can prove the following result.


\begin{theorem}
\source{anticyclotomic-iwasawa-rational-elliptic-eisenstein}
\notready%[Howard]
\label{thm:howard-HPKS}
\source{anticyclotomic-iwasawa-rational-elliptic-eisenstein}
Assume $E(K)[p]=0$. Then there exists a Kolyvagin system $\kappa^{\rm Hg}\in\mathbf{KS}(\mathbf{T},\mathcal{F}_\Lambda,\cL_E)$ such that $\kappa_1^{\rm Hg}\in\rH^1_{\Fcal_\Lambda}(K,\mathbf{T})$ is nonzero.
\end{theorem}

\begin{proof} 
Under the additional hypotheses that $p\nmid h_K$, the class number of $K$, and the representation $G_K\rightarrow{\rm Aut}_{\Z_p}(T)$ is surjective, this is \cite[Thm.~2.3.1]{howard}. In the following paragraphs, we explain how to adapt Howard's arguments to our situation. 

We begin by briefly recalling the construction of $\kappa^{\rm Hg}$ in \cite[\S{2.3}]{howard}. Let $K_k$ be the subfield of $K_\infty$ with $[K_k\colon K]=p^k$. For each $n\in\cN$ set
\[
P_k[n]:={\rm Norm}_{K[np^{d(k)}]/K_k[n]}(P[np^{d(k)}])\in E(K_k[n]),
\]
where $d(k)=\min\{d\in\Z_{\geqslant 0}\colon K_k\subset K[p^{d(k)}]\}$, and $K_k[n]$ denotes the compositum of $K_k$ and $K[n]$. Letting $H_k[n]\subset E(K_k[n])\otimes\Z_p$ be the $\Z_p[{\rm Gal}(K_k[n]/K)]$-submodule generated by $P[n]$ and $P_j[n]$ for $j\leqslant k$, it follows from the Heegner point norm relations \cite[\S{3.1}]{perrinriou} that one can form the $\mathcal{G}(n)$-module
\[
\mathbf{H}[n]:=\varprojlim_k H_k[n].
\]
By \cite[Lem.~2.3.3]{howard}, there is a family 
\[
\{Q[n]=\varprojlim_k Q_k[n]\in\mathbf{H}[n]\}_{n\in\cN}
\]
such that
\begin{equation}\label{eq:multiplier}
\source{anticyclotomic-iwasawa-rational-elliptic-eisenstein}
Q_0[n]=\Phi P[n],\quad\textrm{where}\;
\Phi=
\left\{
\begin{array}{ll}
(p-a_p\sigma_p+\sigma_p^2)(p-a_p\sigma_p^*+\sigma_p^{*2}) & \textrm{if $p$ splits in $K$,}\\[0.2em]
(p+1)^2-a_p^2 & \textrm{if $p$ is inert in $K$,}
\end{array}
\right.
\end{equation}
with $\sigma_p$ and $\sigma_p^*$ the Frobenius elements at the primes above $p$ in the split case, and 
\[
{\rm Norm}_{K_{\infty}[n\ell]/K_{\infty}[n]}Q[n\ell] = a_{\ell}Q[n] 
\]
for all $n\ell\in\cN$. Letting $D_n\in\Z_p[G(n)]$ be Kolyvagin's derivative operators, and choosing a set $S$ of representatives for $\mathcal{G}(n)/G(n)$, the class $\kappa_n\in \rH^1(K,\mathbf{T}/I_n\mathbf{T})$ is defined as the natural image of 
\begin{equation}\label{eq:Qn}
\source{anticyclotomic-iwasawa-rational-elliptic-eisenstein}
\tilde{\kappa}_n:=\sum_{s\in S}sD_nQ[n]\in\mathbf{H}[n]
\end{equation}
under the composite map
\[
\bigl(\mathbf{H}[n]/I_n\mathbf{H}[n]\bigr)^{\mathcal{G}(n)}\xrightarrow{\delta(n)}\rH^1(K[n],\mathbf{T}/I_n\mathbf{T})^{\mathcal{G}(n)}\overset{\simeq}\longleftarrow\rH^1(K,\mathbf{T}/I_n\mathbf{T}),
\]
where $\delta(n)$ is induced by the limit of Kummer maps $\delta_k(n):E(K_k[n])\otimes\bZ_p\rightarrow{\rm H}^1(K_k[n],T)$, and the second arrow is given by restriction. (In our case, that the latter is an isomorphism follows from the fact that the extensions $K[n]$ and $\Q(E[p])$ are linearly disjoint, and $E(K_\infty)[p]=0$.) 

The proof that the classes  $\kappa_n$ land in $\rH^1_{\Fcal_\Lambda}(K,\mathbf{T})$ and can be modified to a system $\kappa^{\rm Hg}=\{\kappa_n^{\rm Hg}\}_{n\in\cN}$ satisfying the Kolyvagin system relations is the same as in \cite[Lem.~2.3.4]{howard} \emph{et seq.}, noting that the arguments proving Lemma~2.3.4 (in the case $v\vert p$) apply almost \emph{verbatim} in the case when $p$ divides the class number of $K$. Finally, that $\kappa_1^{\rm Hg}$ is nonzero follows from the works of Cornut and Vatsal \cite{cornut, vatsal}.
\end{proof}

Applying Theorem~\ref{thm:howard} and Corollary~\ref{cor:howard} to the Kolyvagin system $\kappa^{\rm Hg}$ of Theorem~\ref{thm:howard-HPKS}, we thus obtain the following.

\begin{theorem}
\source{anticyclotomic-iwasawa-rational-elliptic-eisenstein}
\notready\label{thm:howard-HP}
\source{anticyclotomic-iwasawa-rational-elliptic-eisenstein}
Assume $E(K)[p]=0$. Then the module $\rH^1_{\Fcal_\Lambda}(K,\mathbf{T})$ has $\Lambda$-rank one, and there is a finitely generated torsion $\Lambda$-module $M$ such that
\begin{itemize}
\item[(i)] $\mathcal{X}\sim\Lambda\oplus M\oplus M$,
\item[(ii)] ${\rm char}_\Lambda(M)$ divides ${\rm char}_\Lambda\bigl(\rH^1_{\Fcal_\Lambda}(K,\mathbf{T})/\Lambda\kappa_1^{\rm Hg}\bigr)$ in $\Lambda[1/p,1/(\gamma-1)]$.
\end{itemize}
Moreover, if $\rH^1_{\CF}(K,E[p^\infty])$ has $\Z_p$-corank one,  then ${\rm char}_\Lambda(M)$ divides ${\rm char}_\Lambda\bigl(\rH^1_{\Fcal_\Lambda}(K,\mathbf{T})/\Lambda\kappa_1^{\rm Hg}\bigr)$ in $\Lambda[1/p]$.
\end{theorem}

\begin{remark}
\source{anticyclotomic-iwasawa-rational-elliptic-eisenstein}
\notready\label{rem:comparison}
\source{anticyclotomic-iwasawa-rational-elliptic-eisenstein}
For our later use, we compare the class $\kappa_1^{\rm Hg}\in\rH^1_{\Fcal_\Lambda}(K,\mathbf{T})$ from Theorem~\ref{thm:howard-HPKS} with the $\Lambda$-adic class constructed in \cite[\S{5.2}]{cas-hsieh1} (taking for $f$ the newform associated with $E$). 

Denote by $\alpha$ the $p$-adic unit root of $x^2-a_px+p$. With the notations introduced in the proof of Theorem~\ref{thm:howard-HPKS}, define the \emph{$\alpha$-stabilized Heegner point} $P[p^k]_\alpha\in E(K[p^k])\otimes\Z_p$ by 
\begin{equation}\label{eq:stabilized-1}
\source{anticyclotomic-iwasawa-rational-elliptic-eisenstein}
P[p^k]_\alpha:=\left\{
\begin{array}{ll}
P[p^k]-\alpha^{-1}P[p^{k-1}]&\textrm{if $k\geqslant 1$,}\\[0.2em]
u_K^{-1}\bigl(1-\alpha^{-1}\sigma_p\bigr)\bigl(1-\alpha^{-1}\sigma_p^*\bigr)P[1]&\textrm{if $k=0$ and $p$ splits in $K$,}\\
u_K^{-1}\bigl(1-\alpha^{-2}\bigr)P[1]&\textrm{if $k=0$ and $p$ is inert in $K$.}
\end{array}
\right.
\end{equation}

Using the Heegner point norm relations, a straightforward calculation shows that the points $\alpha^{-k}P[p^k]_\alpha$ are norm-compatible.  Letting $\delta:E(K_k)\otimes\Z_p\rightarrow\rH^1(K_k,T_pE)$ be the Kummer map, we may therefore set
\[
\kappa_\infty:=\varprojlim_k\delta(\kappa_k)\in\varprojlim_k\rH^1(K_k,T_pE)\simeq\rH^1(K,\mathbf{T}),
\]
where $\kappa_k=\alpha^{-d(k)}{\rm Norm}_{K[p^{d(k)}]/K_k}(P[p^{d(k)}]_\alpha)$. The inclusion $\kappa_\infty\in\rH^1_{\Fcal_\Lambda}(K,\mathbf{T})$ follows immediately from the construction. 
For the comparison with $\kappa_1^{\rm Hg}$, note that by (\ref{eq:Qn}) the projection ${\rm pr}_K({\kappa}_1^{\rm Hg})$ of $\kappa_1^{\rm Hg}$ to $\rH^1(K,T_pE)$ is given by the Kummer image of ${\rm Norm}_{K[1]/K}(Q_0[1])$, while $\kappa_0$ is the Kummer image of ${\rm Norm}_{K[1]/K}(P[1]_\alpha)$. Thus comparing $(\ref{eq:multiplier})$ and $(\ref{eq:stabilized-1})$ we see that 
\begin{equation}\label{eq:comparison}
\source{anticyclotomic-iwasawa-rational-elliptic-eisenstein}
{\rm pr}_K(\kappa_1^{\rm Hg})=
\begin{cases}
u_K\alpha^2(\beta-1)^2\cdot\kappa_0&\textrm{if $p$ splits in $K$},\\
u_K\alpha^2(\beta^2-1)\cdot\kappa_0&\textrm{if $p$ is inert in $K$,}
\end{cases}
\end{equation}
where $\beta=p\alpha^{-1}$. 
%Since $u_K=1$ by hypothesis (\ref{eq:intro-disc}), 
In particular, $\kappa_\infty$ and $\kappa_1^{\rm Hg}$ generate the same $\Lambda$-submodule of $\rH^1_{\Fcal_\Lambda}(K,\mathbf{T})$. 
\end{remark}



\subsection{Proof of the Iwasawa main conjectures}

Let $\kappa_\infty\in\rH^1_{\Fcal_\Lambda}(K,\mathbf{T})$ be the $\Lambda$-adic Heegner class introduced in Remark~\ref{rem:comparison}, and %as in the Introduction, 
put $\Lambda_{\rm ac}=\Lambda\otimes_{\Z_p}\Q_p$.
{Let ${\rm H}^1_{\Fcal_{\rm Gr}}(K,\mathbf{T})$ be defined just as ${\rm H}^1_{\Fcal_{\rm Gr}}(K,M_E)$ but with 
$\mathbf{T}$ replacing $M_E$ and the conditions on $v$ and $\bar v$ switched.

\begin{prop}
\source{anticyclotomic-iwasawa-rational-elliptic-eisenstein}
\notready\label{prop:equiv-imc}
\source{anticyclotomic-iwasawa-rational-elliptic-eisenstein}
Assume that $p=v\bar{v}$ splits in $K$ and that $E(K)[p]=0$. 
%and let $\Lambda'$ denote either $\Lambda$ or $\Lambda[1/p]$. 
Then the following statements are equivalent:
\begin{enumerate}
\item[(i)] Both ${\rm H}^1_{\Fcal_\Lambda}(K,\mathbf{T})$ and 
$\mathcal{X} = {\rm H}^1_{\Fcal_\Lambda}(K,M_E)^\vee$ have $\Lambda$-rank one, and the divisibility 
\[
%\mathrm{char}_\Lambda({\rm H}^1_{\Fcal_\Lambda}(K,M_E)^\vee_{\rm tors})
\mathrm{char}_\Lambda(\mathcal{X}_{\rm tors}) \supset\mathrm{char}_\Lambda\bigl({\rm H}^1_{\Fcal_\Lambda}(K,\mathbf{T})/\Lambda\kappa_\infty\bigr)^2
\]
holds in $\Lambda_{\rm ac}$.
\item[(ii)]  Both ${\rm H}^1_{\Fcal_{\rm Gr}}(K,\mathbf{T})$ and $\mathfrak{X}_{E} = {\rm H}^1_{\Fcal_{\rm Gr}}(K,M_E)^\vee$ are $\Lambda$-torsion, and the divisibility 
\[
\mathrm{char}_\Lambda(\mathfrak{X}_{E})\Lambda^{\rm ur}\supset(\mathcal{L}_E)
\]
holds in $\Lambda^{\rm ur}\otimes_{\bZ_p}\bQ_p$.
%_{\rm ac}\hat{\otimes}_{\bZ_p}\bZ_p^{\rm ur}$
\end{enumerate}
Moreover, the same result holds for the opposite divisibilities.
\end{prop}
}

\begin{proof}
See \cite[Thm.~5.2]{BCK}, whose proof still applies after inverting $p$. 
\end{proof}

We can now conclude the proof of Theorem~\ref{thm:C} in the introduction. 

\begin{theorem}
\source{anticyclotomic-iwasawa-rational-elliptic-eisenstein}
\notready\label{thm:thmA}
\source{anticyclotomic-iwasawa-rational-elliptic-eisenstein}
Suppose $K$ satisfies hypotheses {\rm (\ref{eq:intro-Heeg})}, {\rm (\ref{eq:intro-spl})}, {\rm (\ref{eq:intro-disc})}, and 
{\rm (\ref{eq:intro-sel})}, and that
$E[p]^{ss}=\mathbb{F}_p(\phi)\oplus\mathbb{F}_p(\psi)$ as $G_\bQ$-modules, with $\phi\vert_{G_p}\neq\mathds{1},\omega$.
Then $\mathfrak{X}_E$ is $\Lambda$-torsion, and
\[
{\rm char}_\Lambda(\mathfrak{X}_E)\Lambda^{\rm ur}=(\mathcal{L}_E)
\]
as ideals in $\Lambda^{\rm ur}$. 
\end{theorem}

\begin{proof}
By Theorem~\ref{thm:howard-HP}, 
%we can appeal to Theorem~\ref{thm:howard} and especially Corollary~\ref{cor:howard} with $\kappa=\kappa^{\rm Hg}$, and so
the modules ${\rm H}^1_{\Fcal_\Lambda}(K,\mathbf{T})$ and 
$\rH^1_{\Fcal_\Lambda}(K,M_E)^\vee$ have both $\Lambda$-rank one, with
\[
{\rm char}_\Lambda\bigl(\rH^1_{\Fcal_\Lambda}(K,M_E)^\vee_{\rm tors}\bigr)\supset{\rm char}_\Lambda\bigl({\rm H}^1_{\Fcal_\Lambda}(K,\mathbf{T})/\Lambda\kappa_1^{\rm Hg}\bigr)^2
\]
as ideals in $\Lambda_{\rm ac}=\Lambda[1/p]$. Since by Remark~\ref{rem:comparison} 
the classes $\kappa_1^{\rm Hg}$ and $\kappa_\infty$ generate the same $\Lambda$-submodule of $\rH^1_{\Fcal_\Lambda}(K,\mathbf{T})$, by Proposition~\ref{prop:equiv-imc} it follows that $\mathfrak{X}_E$ is $\Lambda$-torsion, with
\[
{\rm char}_\Lambda(\mathfrak{X}_E)\Lambda^{\rm ur}\supset(\mathcal{L}_E)
\]
as ideals in $\Lambda_{\rm ac}\hat{\otimes}_{\bZ_p}\bZ_p^{\rm ur}$. This divisibility, together with the equalities $\mu(\mathfrak{X}_E)=\mu(\mathcal{L}_E)=0$ and $\lambda(\mathfrak{X}_E)=\lambda(\mathcal{L}_E)$ in Theorem~\ref{mulambda}, yields the result.
\end{proof}

As a consequence, we can also deduce the first cases of Perrin-Riou's Heegner point main conjecture \cite{perrinriou} in the residually reducible case. More precisely, together with Theorem~\ref{thm:howard-HP}, the following yields Corollary~\ref{cor:D} in the introduction. 

\begin{cor}
\source{anticyclotomic-iwasawa-rational-elliptic-eisenstein}
\notready\label{cor:PR} Suppose $K$ satisfies hypotheses {\rm (\ref{eq:intro-Heeg})}, {\rm (\ref{eq:intro-spl})},  {\rm (\ref{eq:intro-disc})}, and {\rm (\ref{eq:intro-sel})}, and that
\source{anticyclotomic-iwasawa-rational-elliptic-eisenstein}
$E[p]^{ss}=\mathbb{F}_p(\phi)\oplus\mathbb{F}_p(\psi)$ as $G_\bQ$-modules, with $\phi\vert_{G_p}\neq\mathds{1},\omega$. 
Then both ${\rm H}^1_{\Fcal_\Lambda}(K,\mathbf{T})$ and 
$\rH^1_{\Fcal_\Lambda}(K,M_E)^\vee$ have $\Lambda$-rank one, and
\[
\mathrm{char}_\Lambda(\rH^1_{\Fcal_\Lambda}(K,M_E)^\vee_{\rm tors})=\mathrm{char}_\Lambda\bigl({\rm H}^1_{\Fcal_\Lambda}(K,\mathbf{T})/\Lambda\kappa_\infty\bigr)^2
\]
as ideals in $\Lambda_{\rm ac}$.
\end{cor}

\begin{proof}
In light of Remark~\ref{rem:comparison}, this is the combination of Theorem~\ref{thm:thmA} and Proposition~\ref{prop:equiv-imc}.
%Everything except one of the divisibilities of the characteristic ideals has been shown in Theorem~\ref{thm:howard}; the missing divisibility now follows from Theorem~\ref{thm:thmA} together with Proposition~\ref{prop:equiv-imc}.
\end{proof}

\begin{remark}
\source{anticyclotomic-iwasawa-rational-elliptic-eisenstein}
\notready\label{rmk:hyp-sel}
\source{anticyclotomic-iwasawa-rational-elliptic-eisenstein}
If the Heeger point $P_K = {\rm Norm}_{K[1]/K}(P[1])\in E(K)$ is non-torsion, then
{\rm (\ref{eq:intro-sel})} holds by the main results of \cite{kolyvagin-mw-sha}.
In particular, this is so if the image ${\rm pr}_K(\kappa_1^{\rm Hg})$ of $\kappa_1^{\rm Hg}$ 
(equivalently, the class $\kappa_0$) in $\rH^1_{\CF}(K,T_p(E))$ is non-zero 
(as $\rH^1_{\CF}(K,T_p(E))$ is non-torsion since $E(K)[p]= 0$). 
\end{remark}
